\section{Introduction}
\label{sec:introduction}

Streaming machine learning-based intrusion detection systems (IDS) increasingly execute at the network edge to eliminate raw traffic egress bandwidth, minimize alert latency, and preserve data privacy~\cite{sforzin2016rpids,baalaji2025ensemble}. Edge deployments on embedded System-on-Chip (SoC) accelerators---such as the NVIDIA Jetson platform---routinely couple a high-performance primary serving runtime (\eg ONNX Runtime leveraging C++ SIMD and GPU tensor cores) with an in-process, zero-dependency fallback engine (\eg written in pure NumPy/C) designed to handle cold-start initialization, container failover, and memory-constrained recovery~\cite{asghar2026hedid,hasan2025iiot}.

However, providing dual-runtime serving for mission-critical security applications introduces a fundamental correctness challenge: \textbf{Runtime Equivalence}. A detector that serves identical artifacts across two distinct runtimes promises the exact same intrusion verdict from both. Yet, in practice, disparate runtimes implement different internal floating-point precisions (\eg float32 vs.\ float64), diverse vectorization instruction sets, and divergent tie-breaking semantics at decision thresholds~\cite{jajal2024onnxconverters,wen2025mindthegap}. If a high-volume flow falls on a split boundary where runtime arithmetic diverges, verdicts flip silently between engines.

Standard validation protocols in literature rely almost exclusively on replaying recorded network traffic traces (\eg replaying held-out test splits from CICIDS2017 or CSE-CIC-IDS2018). We observe that this practice induces a severe false sense of security: observed network flows almost never coincide with exact floating-point decision thresholds. In empirical evaluations, millions of replayed flows fail to stimulate the critical boundary regions where arithmetic discrepancies manifest, allowing severe runtime divergences to escape into production unnoticed.

To overcome these fundamental limitations, this paper establishes a mathematically formal and empirically validated methodology for high-throughput, parity-verified tree-ensemble serving at the network edge. Specifically, we make four principal contributions:

\begin{enumerate}
    \item \textbf{Boundary-Directed Parity Testing:} We formalize Algorithm~\ref{alg:boundary_probe}, an artifact-driven probe generation algorithm that operates without requiring real-world traffic data. By climbing ancestor paths in the tree topology to compute feasible bounding hyper-rectangles, the algorithm synthesizes exact flow vectors that reach internal split nodes and systematically probes their $\pm 1\,\text{ulp}$ floating-point neighborhoods. On deployed production forests, our method exposes precision discrepancies $34\times$ more densely than replaying complete test splits, proving that standard $2{,}000$-row replay gates miss rounding divergences over $93.5\%$ of the time.
    \item \textbf{Cache-Blocked Lockstep Traversal with Exact Early Exit:} To eliminate the debilitating interpreter overhead of naive sequential tree traversals, we design Algorithm~\ref{alg:lockstep_early_exit}. By structuring tree evaluation into cache-sized blocks ($\le 8192$ sample-tree pairs) that fit within L1/L2 caches, the algorithm achieves lockstep SIMD execution across multiple trees. Furthermore, by tracking remaining extremal margin bounds $[R^{\min}, R^{\max}]$, we establish an exact early exit rule that provably guarantees zero label divergence ($0\,\text{flip}$) while dramatically reducing average tree evaluation depth.
    \item \textbf{Empirical Edge Validation on Jetson Orin Nano Testbed:} We implement and deploy our complete serving pipeline across two NVIDIA Jetson Orin Nano Super boards linked via a brokerless, zero-copy ZeroMQ streaming bus. We eliminate JVM broker dependencies and serialization bottlenecks, demonstrating in-process serving throughput of up to $21{,}191\,\text{flows/s}$ ($0.057\,\text{ms}$ latency) on physical edge hardware and proving that machine learning inference is completely removed from the critical latency path.
    \item \textbf{Cross-Paradigm and Cross-Dataset Generalization:} We demonstrate that boundary threshold sensitivity is not an artifact of a single dataset. Across four distinct training paradigms (RandomForest, HistGradientBoosting, LightGBM, and XGBoost) and multiple benchmark datasets, we prove that histogram binning inherently produces dense threshold sharing where float precision divergence induces up to a $1.58\%$ label flip rate. In contrast, our float32 lockstep traversal maintains exact, bitwise parity across all models.
\end{enumerate}

The remainder of this paper is organized as follows: Section~\ref{sec:related_work} reviews related work in edge IDS and runtime verification. Section~\ref{sec:methodology} introduces the formal methodology, including Algorithm~\ref{alg:boundary_probe} and Algorithm~\ref{alg:lockstep_early_exit}. Section~\ref{sec:architecture} details the edge architecture and testbed deployment. Section~\ref{sec:evaluation} presents comprehensive empirical benchmarks on parity, throughput, and cold start. Section~\ref{sec:generalization} analyzes cross-dataset and cross-framework generalization. Section~\ref{sec:discussion} discusses architectural trade-offs and threats to validity, and Section~\ref{sec:conclusion} concludes the paper.
